% CONDITIONAL SECTION — included under Option A of paper/PLAN.md §1.3 (scope
% ruling still open for T0). Under Option B this file is deleted wholesale,
% together with 06-lattice.tex; no other section depends on it except via the
% label sec:framework.
% ============================================================================
% 08-dolgachev-doran.tex — the M_n-polarized framework. Claims #18-#20 of
% PLAN.md §3. rho=19 / T=3 is presented CONDITIONALLY (PLAN.md §5 item 4); the
% identification of the computed lattice with T is conjectural throughout.
% ============================================================================
\section{Relation to the \texorpdfstring{$M_n$}{Mn}-polarized framework}
\label{sec:framework}

This section places the computations of \S\S\ref{sec:sym2}--\ref{sec:integrality}
next to a classical framework and states carefully what the juxtaposition does
and does not establish. Nothing in this section is claimed as a new theorem.

\subsection{The framework, as cited}

\begin{theorem}[cited, Dolgachev 1996]\label{thm:dolgachev}
\cite[\S7, p.~20 and Thm.~7.1]{Dolgachev1996}. For $M_n = \U\oplus E_8^2
\oplus\lat{-2n}$ one has $(M_n)^{\perp} = \U\oplus\lat{2n}$ in the K3 lattice,
and the coarse moduli space of $M_n$-polarized K3 surfaces is the Fricke
modular curve $\mathbb{H}/\Gamma_0(n)+$.
\end{theorem}

\begin{theorem}[cited, Doran 2000]\label{thm:doran}
\cite[Thm.~5.13]{Doran2000}. The Picard--Fuchs operator of an
$M_n$-polarized family of K3 surfaces is the symmetric square of a
second-order Fuchsian operator.
\end{theorem}

Both sources were fetched, hash-pinned, and read by the program; the parts of
Theorem~\ref{thm:dolgachev} used below were additionally re-derived
symbolically before use, with the period shape and the invariant bilinear form
recomputed from scratch rather than transcribed. We also record, because it
governs how much the coincidences below can be made to prove, that Doran
himself notes the classification of rank-19 lattice polarizations to be
incomplete~\cite[\S6]{Doran2000}.

\subsection{The logical direction}

A referee's first reaction to \S\ref{sec:lattice} should be that
Theorem~\ref{thm:dolgachev} already says $(M_7)^{\perp} = \U\oplus\lat{14}$
and Theorem~\ref{thm:doran} already says the Picard--Fuchs operator of such a
family is a symmetric square, so nothing has been learned. The reply is that
the direction of inference here is the reverse of the framework's. The
framework starts from a polarization hypothesis and deduces invariants. We
start from a concretely given operator family --- specified by four rational
numbers --- and \emph{compute} its invariants: the symmetric-square structure
(kernel-checked, and for the whole template rather than for one member), the
Fuchsian signature, the modular coordinate and its level, and the monodromy
lattice. Each of these is a fact about the operator, established without any
geometric hypothesis. The framework then supplies a reading under which they
all cohere, and the coherence is close: the same level $7$ appears from four
computations that share no code, and the monodromy lattice matches
$(M_7)^{\perp}$ on the nose, including the discriminant form.

That is convergent evidence, not a proof, and the incompleteness of the
rank-19 classification is precisely why it cannot be upgraded by citation.

\subsection{The evidence, assembled}

\begin{table}[htbp]
\centering
\begin{tabular}{@{}p{0.42\textwidth}ll@{}}
\toprule
\textbf{Computed fact about the $s_7$ operators} & \textbf{Status} &
\textbf{Framework counterpart} \\
\midrule
$\Lthree = P_2\cdot\Symtwo(\Ltwo)$, uniformly on the template
  (Thm.~\ref{thm:sym2-template}) & (K) & Thm.~\ref{thm:doran} \\
Almkvist--van Straten $W\equiv0$ (Thm.~\ref{thm:avs}) & (K) &
  self-adjointness of a $\Symtwo$ \\
$\Ltwo$, $\Lthree$ irreducible; $\Lthree$ minimal
  (\S\ref{sec:irreducibility}) & (E) & rank-3 sub-system \\
Implied Fuchsian signature: genus $0$, orders $(2,2,3)$, one cusp,
  area$/2\pi = \tfrac23$ (Rem.~\ref{rem:signature}) & (E) &
  signature of $\Gamma_0(7)+$ \\
Modular coordinate is a $\Gamma_0(7)+$ Hauptmodul, $\kappa = 49$
  (Prop.~\ref{prop:hauptmodul}) & (E) & $\mathbb{H}/\Gamma_0(7)+$ of
  Thm.~\ref{thm:dolgachev} \\
$\sqrt7$ in the elliptic monodromies (NC~\ref{nc:monodromy}) & (N) &
  level $7$ \\
Monodromy lattice $\U\oplus\lat{14}$, $\det = -14$, disc.\ $\Z/14$,
  $q = \tfrac1{14}$ (Props.~\ref{prop:lattice},~\ref{prop:usplit}) & (N)+(E) &
  $(M_7)^{\perp}$ of Thm.~\ref{thm:dolgachev} \\
$s_{10}$ control gives $\U\oplus\lat{20}$ (NC~\ref{nc:s10-control}) & (N)+(E) &
  level $10$, not level $7$ \\
Orthogonal complement of $G$ in $\Lambda$ is
  $\U\oplus E_8(-1)^2\oplus\lat{-14}$ (Prop.~\ref{prop:g0complement}) & (N)+(E) &
  exactly $M_7$, mutual complement of Thm.~\ref{thm:dolgachev} \\
\bottomrule
\end{tabular}
\medskip
\caption{Computed invariants of the $s_7$ operator family against the
$M_n$-polarized framework. Status labels as in \S\ref{sec:intro}.}
\label{tab:evidence}
\end{table}

\subsection{A conditional rank statement}

We now record the rank statement that the minimality result of
\S\ref{sec:irreducibility} yields \emph{if} two external inputs are granted.
It is stated as a conditional proposition, with its hypotheses displayed, and
it is used nowhere else in this paper.

\begin{proposition}[conditional; status (C)+(E) under (H1)--(H2)]
\label{prop:ranks}
Assume:
\begin{itemize}
\item[\textbf{(H1)}] the $s_7$ family is realized by the explicit projective
  K3 model of Almkvist and van Straten~\cite{AlmkvistvanStraten2021}, and
  $\Lthree$ is the Picard--Fuchs operator of the period of its holomorphic
  $2$-form; and
\item[\textbf{(H2)}] for the very general member $X$ of the family,
  $\operatorname{rank}T(X)$ equals the rank of the $\Q$-sub-variation of Hodge
  structure generated by that period under Gauss--Manin.
\end{itemize}
Then $\operatorname{rank}T(X) = 3$ and $\rho(X) = 19$.
\end{proposition}

\begin{proof}
By Corollary~\ref{cor:minimal}, the sub-variation generated by the period has
rank $3$, that being the order of its minimal annihilator. (H2) transfers this
to $\operatorname{rank}T(X)$, and $\rho(X) = 22 - \operatorname{rank}T(X) = 19$.
\end{proof}

\begin{remark}[On (H2), and on the caveat that must travel with $\rho = 19$]
\label{rem:H2}
(H2) is standard Hodge theory and we cite rather than re-derive it. It is the
conjunction of two inclusions. First, for very general $z$ the N\'eron--Severi
lattice is locally constant, the Noether--Lefschetz locus being a countable
union of proper subvarieties, so $T = \NS^{\perp}$ is a Gauss--Manin-flat
sub-Hodge-structure containing the holomorphic $2$-form; the sub-variation
generated by that form is therefore contained in $T$. Second, $T(X)\otimes\Q$
is by definition the minimal primitive sub-Hodge-structure of $H^2(X,\Q)$
whose complexification contains $H^{2,0}$, and the saturation of the
sub-variation is such a structure; so the reverse inclusion holds
\cite[ch.~3 \S3.1, ch.~17 \S17.1]{Huybrechts2016}. Simplicity of the
transcendental Hodge structure is~\cite[Thm.~1.6(a)]{Zarhin1983}.

The caveat is essential and is not a formality: $\rho = 19$ is the value for
the \emph{very general} member. It jumps to $20$ on a countable dense subset
of the base. Any use of Proposition~\ref{prop:ranks} that drops the phrase
``very general'' is a misuse of it.
\end{remark}

\subsection{An independent cross-check: the orthogonal complement (G0)}
\label{subsec:g0}

A separate, purely lattice-theoretic computation supplies additional support
for the coherence recorded in Table~\ref{tab:evidence}, and it is reported
here with exactly the same conditionality as \S\ref{sec:lattice}: it takes the
recognized matrices of Numerical Claim~\ref{nc:monodromy} as input, is
otherwise exact, and is dated 2026-07-28.

\begin{proposition}[The orthogonal complement of $G$ in the K3 lattice;
status (E), given Numerical Claim~\ref{nc:monodromy}]\label{prop:g0complement}
Let $G$ be the Gram matrix of Proposition~\ref{prop:lattice}(iii), embedded
primitively in the K3 lattice $\Lambda = \U^3\oplus E_8(-1)^2$ via the
hyperbolic-plane witness $f\mapsto e_{U_1},\ e\mapsto f_{U_1},\ w\mapsto
e_{U_2}+7f_{U_2}$ (a norm-$14$ vector in a hyperbolic plane is forced to have
this shape; there is no other primitive choice). Its orthogonal complement is
isometric over $\Z$ to
\[
  \U \oplus E_8(-1) \oplus E_8(-1) \oplus \lat{-14},
\]
of rank $19$, signature $(1,18)$, cyclic discriminant group of order $14$
(discriminant-form value $q = \tfrac{27}{14} \bmod 2\Z$ on a generator), and
existence and uniqueness of the primitive embedding are certified rather than
assumed, via the Nikulin/Eichler bound $\operatorname{rank}\Lambda -
\operatorname{rank}G = 19 \ge \ell(A_G) + 2 = 3$.
\end{proposition}

Two remarks fix what this proposition does and does not add to
Proposition~\ref{prop:ranks}.

\emph{What it does not add.} Arithmetically, rank $19$ here is forced by
$\operatorname{rank}\Lambda - \operatorname{rank}G = 22 - 3$; it carries no
information about $\rho$ beyond what Proposition~\ref{prop:ranks} already
gives, and it is \emph{not} presented as a second derivation of $\rho = 19$.
Nor does it discharge hypotheses (H1)/(H2) of Proposition~\ref{prop:ranks}, or
either gap of Conjecture~\ref{conj:T} below: it is a further exact computation
downstream of the same numerically-recognized input (Numerical
Claim~\ref{nc:monodromy}), not an independent route to the identification of
the transcendental lattice.

\emph{What it does add.} The proposition computes the \emph{explicit isometry
type} of the complement, not merely its rank. Comparing it to
Theorem~\ref{thm:dolgachev}: under the standard convention that $E_8^2$
there denotes $E_8(-1)\oplus E_8(-1)$ (as fixed by $\Lambda$ itself, above),
$\U\oplus E_8(-1)^2\oplus\lat{-14}$ is exactly $M_7$. On the computation, $G$
and $M_7$ are thus mutual orthogonal complements inside $\Lambda$ --- the same
coherence already recorded for $T$ and $(M_7)^{\perp}$ in
Table~\ref{tab:evidence}, now confirmed from the other side, together with a
verified (not assumed) uniqueness-of-embedding bound.

On independence, precisely: this computation has been carried out three
times --- two in-house exact-symbolic implementations sharing no code (a
constructive embedding witness, and an independent genus-uniqueness argument),
agreeing on the Gram matrix bit-for-bit, and one further zero-shot
re-derivation by a separate reviewing model, which built its own $E_8$ Cartan
matrix from a different Dynkin-node labeling and independently recomputed the
discriminant value; all three agree exactly
\cite{G0NSGenus2026,G0DecisionLog2026}. This is reported as corroboration of
the computation, not as a fourth independent geometric derivation of $\rho=19$:
the third check was supplied the specific embedding formula displayed above
rather than asked to discover it, and, like the first two, it starts from the
same Numerical Claim~\ref{nc:monodromy} input.

\subsection{The lattice arithmetic, kernel-checked}\label{subsec:mn-lean}

The identification of the computed lattice with $\U\oplus\lat{14}$ is status
(E) and is not re-derived in Lean. What \emph{is} kernel-checked
(\artifact{Agora/Geometry/MnLattice.lean}) is the lattice arithmetic that the
identification is then fed into, uniformly in $N$. It is built on the lattice
interface of an independent Lean development of the K3 and Narain
lattices~\cite{LeanMaster2026}, which supplies the Gram-matrix type, the
hyperbolic plane, and the signature of $\Lambda_{K3}$; nothing in the present
subsection depends on that development's physical content.

\begin{proposition}[status (K)]\label{prop:mn-lean}
Let $N\in\Z$ and let $T_N$ be the Gram matrix of $\U\oplus\lat{2N}$ in a basis
$(e,f,w)$ with $e^2=f^2=0$, $e\cdot f=1$, $w^2=2N$.
\begin{enumerate}
\item $T_N$ is symmetric and even, with $\det T_N=-2N$; for $N=7$ it is not
  unimodular.
\item \emph{(Glue.)} In a single hyperbolic plane, $e+Nf$ and $e-Nf$ are
  orthogonal of norms $2N$ and $-2N$, and span a sublattice of index $2N$;
  $e+Nf$ is primitive. For $N=7$ this is the witness $w\mapsto e+7f$ of
  Proposition~\ref{prop:g0complement}.
\item \emph{(Swap.)} The exchange $\sigma\colon e\leftrightarrow f$ is an
  isometry of $T_N$, an involution, of determinant $-1$.
\item \emph{(Swap $=$ Fricke.)} Over any field, the vector
  $\omega(\tau)=e-N\tau^2f+\tau w$ is isotropic, and for $N,\tau\neq0$
  \[
    \sigma\,\omega(\tau) \;=\; (-N\tau^2)\;\omega\!\Bigl(-\frac{1}{N\tau}\Bigr).
  \]
  If $N\tau^2=-1$ then $\sigma\,\omega(\tau)=\omega(\tau)$.
\end{enumerate}
\end{proposition}

Item~(iv) is the statement, at the level of the period line, that the lattice
involution $\sigma$ induces the Fricke involution $\tau\mapsto-1/(N\tau)$ on
the tube-domain parameter of the $M_N$-polarized period domain
\cite{Dolgachev1996}; the formalization includes a control refuting the
variant with the opposite sign, so the statement is not an artifact of the
normalization. The signatures $(2,1)$ and $(1,18)$ assigned to
$\U\oplus\lat{2N}$ and $M_N$ add to the $(3,19)$ of $\Lambda_{K3}$
(\artifact{sig\_fills\_K3}); as in the cited development, these are asserted
pairs, not extracted from definiteness proofs, and the rank $19$ is forced by
$22-3$.

One structural observation, stated with its limits. The $2\times2$ block of
$\sigma$ on the $\U$ summand is the Gram matrix of $\U$ itself, which the cited
development proves equal to the Narain form of a circle compactification, where
the same exchange (of momentum and winding) is the generator of T-duality. So a
single integer matrix is, on one lattice, the Fricke involution of the $M_N$
period domain and, on another, the T-duality generator. This is a fact about a
lattice isometry. We do not assert, and nothing here supports, a physical
identification of the two.

\subsection{The modular group as a group of lattice isometries}\label{subsec:modular-action}

Proposition~\ref{prop:mn-lean}(iii)--(iv) exhibited one isometry of
$\U\oplus\lat{2N}$ acting on the period as the Fricke involution. That was a
sample of a general structure, which we now state in the form in which it is
kernel-checked (\artifact{Agora/Geometry/ModularAction.lean}). Write
$\omega(\tau) = e - N\tau^2 f + \tau w$ as before.

\begin{proposition}[status (K)]\label{prop:rho}
Over any commutative ring, define
\[
  \rho(a,b,c,d) \;=\;
  \begin{pmatrix}
    d^2 & -Nc^2 & 2Ncd\\
    -Nb^2 & a^2 & -2Nab\\
    bd & -ac & ad+Nbc
  \end{pmatrix},
  \qquad
  \rho_W(a,b,c,d) \;=\;
  \begin{pmatrix}
    Nd^2 & -c^2 & 2Ncd\\
    -b^2 & Na^2 & -2Nab\\
    bd & -ac & Nad+bc
  \end{pmatrix}.
\]
Then:
\begin{enumerate}
\item $\rho(a,b,c,d)$ is an isometry of $\U\oplus\lat{2N}$ whenever
  $(ad-Nbc)^2=1$, and $\rho_W(a,b,c,d)$ is one whenever $(Nad-bc)^2=1$.
\item $\rho$ is multiplicative: $\rho(g)\rho(g') = \rho(gg')$, the product being
  that of $\Gamma_0(N)$ written on the entries.
\item \emph{(Automorphy.)} As an identity of polynomials --- no determinant
  hypothesis, no invertibility ---
  \[
    \rho(a,b,c,d)\,\omega(\tau)
      \;=\; \bigl((Nc\tau+d)^2,\; -N(a\tau+b)^2,\; (a\tau+b)(Nc\tau+d)\bigr),
  \]
  which over a field is $(Nc\tau+d)^2\,\omega\!\bigl(\tfrac{a\tau+b}{Nc\tau+d}\bigr)$.
  Similarly $\rho_W$ induces $\tau\mapsto\frac{Na\tau+b}{N(c\tau+d)}$.
\end{enumerate}
\end{proposition}

So $\rho$ realizes $\Gamma_0(N)$, and $\rho_W$ the Atkin--Lehner elements
$W=\left(\begin{smallmatrix}Na&b\\Nc&Nd\end{smallmatrix}\right)/\sqrt N$, inside
$O(\U\oplus\lat{2N})$, acting on the period line with the weight-2 automorphy
factor; the $\sqrt N$ cancels in the symmetric square, which is why
Atkin--Lehner involutions are lattice isometries at all. This is the
constructive half of Dolgachev's $O^{+}(\U\oplus\lat{2n})/\pm1\cong\Gamma_0(n)^{+}$
\cite{Dolgachev1996}; the converse is quoted, not proved. The Fricke involution
is $\rho_W(0,-1,1,0)$, which is $-\sigma$ for the $\sigma$ of
Proposition~\ref{prop:mn-lean}, so that proposition is the specialization of this
one. A control confirms the determinant hypothesis in (i) is load-bearing: at
$(a,b,c,d)=(1,0,0,2)$, $N=7$, the matrix is not an isometry.

\subsection{The singular points of the operator are self-dual points}\label{subsec:selfdual}

We can now say what the two finite singular points of $\Lthree$ are. Recall from
\S\ref{sec:operators} that the leading coefficient of the $s_7$ partner operator
is $P_2 = 1-26z-27z^2$, whose roots are $z=-1$ and $z=\tfrac1{27}$; an earlier
reading of these as Kodaira degenerations was withdrawn, and they were
subsequently identified as order-2 elliptic points of $X_0(7)^{+}$.

\begin{proposition}[status (K)]\label{prop:selfdual}
Let $z(h) = h/(1+13h+49h^2)$. Then, over any field:
\begin{enumerate}
\item $z\bigl(\tfrac1{49h}\bigr) = z(h)$, so $z$ is invariant under the Fricke
  involution of the $h$-line;
\item over $\Q$, the fixed points of $h\mapsto \tfrac1{49h}$ are exactly
  $h=\pm\tfrac17$, and $z(\tfrac17)=\tfrac1{27}$, $z(-\tfrac17)=-1$;
\item $\bigl(1-26z(h)-27z(h)^2\bigr)\,(1+13h+49h^2)^2 = (1-49h^2)^2$.
\end{enumerate}
Moreover $\langle e-f,\,\omega(\tau)\rangle = -(N\tau^2+1)$, so the locus
$N\tau^2=-1$ fixed by the Fricke involution is exactly the wall of the
$(-2)$-root $e-f$; and for $N=7$ the conjugate Atkin--Lehner element
$\rho_W(1,4,-2,-1)$ is an involution whose negative is the reflection in the
$(-2)$-root $(-2,4,1)$ of $\U\oplus\lat{14}$.
\end{proposition}

Item (iii) is the sharp statement: pulled back to the $h$-line, the leading
coefficient of the operator becomes a \emph{perfect square} vanishing exactly at
the two Fricke fixed points. The singular locus of $\Ltwo,\Lthree$ is the
ramification locus of the quotient by the involution, and both of its points are
walls of $(-2)$-roots of the lattice --- one for the Fricke involution, one for
its conjugate. We regard this as the correct replacement for the withdrawn
Kodaira reading: it is an exact statement, it is kernel-checked, and it explains
why there are exactly two such points and why they are the ones they are.

\begin{remark}[What the identification costs]\label{rem:selfdual-status}
Proposition~\ref{prop:selfdual} is a theorem about the rational map $z(h)$ and
about the lattice. Reading it as a statement about the $s_7$ \emph{family}
requires two further inputs, neither kernel-checked. First, that
$h=\bigl(\eta(7\tau)/\eta(\tau)\bigr)^4$ and $z=z(h)$ pull the family back to the
$\tau$-line: we verify the equivalent $q$-series identity
$\sum_n s_7(n)z(q)^n = \bigl(\sum_{a,b}q^{a^2+ab+2b^2}\bigr)^2$ exactly to
$O(q^{40})$, status (E), with a control in which replacing $13$ by $12$ makes the
identity fail. Second, that $h(-1/(7\tau)) = 1/(49h(\tau))$, which is the
$\eta$~transformation law, status (L); numerically $h(i/\sqrt7)=\tfrac17$ and
$h\bigl((-1+i/\sqrt7)/2\bigr)=-\tfrac17$ to 40 digits, and the Fricke relation
holds to $10^{-42}$ at a generic point. The script is
\artifact{scripts/check\_selfdual\_points\_s7.py}.
\end{remark}

\begin{remark}[No physical claim]\label{rem:selfdual-nophysics}
In the string-theory literature the wall of a $(-2)$-root is associated with
enhanced gauge symmetry, and the fixed locus of T-duality on a circle is the
self-dual radius; the $2\times2$ block of $\sigma$ on the $\U$ summand is
literally the Narain Gram matrix of a circle, on which the same exchange is the
T-duality generator. We record this because it is what makes the structure
interesting, and we claim nothing from it: no compactification is constructed
here, and a coincidence of lattice isometries is not a physical identification.
\end{remark}

\subsection{What is conjectural}

\begin{conjecture}[Identification of the computed lattice]\label{conj:T}
The joint monodromy-invariant lattice computed in \S\ref{sec:lattice} is the
transcendental lattice of the very general member of the $s_7$ family; that
is, $T \cong \U\oplus\lat{14}$, and the family is $M_7$-polarized.
\end{conjecture}

We are explicit about the two gaps between Proposition~\ref{prop:lattice} and
Conjecture~\ref{conj:T}, because within the computation itself the
identification is pinned only up to two enumerable ambiguities, and only one
of them is closed by computation.

\begin{enumerate}
\item \emph{Even invariant overlattices.} A priori the transcendental lattice
  could be an even overlattice of the computed orbit lattice, with the same
  rational form. These were enumerated: there are none
  (Proposition~\ref{prop:lattice}(v)). This gap is closed, conditionally on
  Numerical Claim~\ref{nc:monodromy}.
\item \emph{Overall integral rescaling.} The transcendental lattice could
  carry $\lambda\cdot G$ for an integer $\lambda>1$. This branch is excluded
  by the shape $\U\oplus\lat{2n}$ that the framework predicts
  (Theorem~\ref{thm:dolgachev}) --- that is, by a reading, not by a
  computation. It is the honest residual gap, and we do not paper over it: the
  argument that rules it out presupposes the conclusion it is used to support.
\end{enumerate}

To these must be added the fact that Numerical Claim~\ref{nc:monodromy} is
itself numerics-backed, and hypothesis (H1) of
Proposition~\ref{prop:ranks}, which is a citation about a geometric model and
is not verified anywhere in this paper. Proposition~\ref{prop:g0complement}
(\S\ref{subsec:g0}) closes neither gap: it computes the orthogonal complement
of the same conjectured lattice, and so is exactly as conditional as the
lattice itself.

\subsection{What is not claimed}

We close by listing, explicitly, statements a reader might expect to find here
and will not.

\begin{itemize}
\item No degeneration or fibre-type reading is made of the singular points
  $z=-1$ and $z=\tfrac1{27}$. Remark~\ref{rem:signature} identifies them as
  order-2 elliptic points of the quotient curve; that is the only reading
  taken, and no elliptic fibration is invoked anywhere.
\item No claim is made that the $s_{10}$ or $s_{18}$ partners are modular,
  and no level is assigned to $s_{18}$ (Remark~\ref{rem:integrality-scope}).
\item Theorem~\ref{thm:sym2-template} is not claimed to imply
  Theorem~\ref{thm:doran} or to be implied by it: the template theorem
  exhibits a partner from four rational parameters with no geometric
  hypothesis, whereas Doran's theorem asserts existence under a polarization
  hypothesis. Neither subsumes the other.
\item Proposition~\ref{prop:ranks} is not asserted unconditionally, and
  Conjecture~\ref{conj:T} is not asserted at all.
\item Proposition~\ref{prop:g0complement} (\S\ref{subsec:g0}) is not claimed
  as evidence independent of \S\ref{sec:lattice}: it shares Numerical
  Claim~\ref{nc:monodromy} with every other statement in this section, and its
  three-way agreement is a corroboration of one exact computation, not three
  independent geometric derivations of $\rho=19$.
\item \textbf{[STUB, 2026-07-30]} No claim of the form ``$T(\text{$s_{10}$}) \cong
  \U\oplus\lat{20}$'' is made anywhere in this paper, and none is available to
  make: the only extant occurrence of $\U\oplus\lat{20}$ in the program is
  Numerical Claim~\ref{nc:s10-control} of \S\ref{sec:lattice}, a
  \emph{discriminating control} run on the identical pipeline to show
  level-sensitivity against the $s_7$ result (Remark~\ref{rem:lattice-scope}
  states explicitly that even the $s_7$ lattice's identification with a
  transcendental lattice is not established there, let alone $s_{10}$'s). It
  is not a certified, independently-verified monodromy-lattice derivation for
  $s_{10}$ analogous to Proposition~\ref{prop:usplit}/\ref{rem:two-verifications}
  for $s_7$: no second implementation, no negative-control suite, and no
  $s_{10}$-specific G0-style orthogonal-complement check exist for it in this
  repository or in Stream~2's epistemic ledger (\texttt{PREDICTION.md},
  \texttt{PROOF\_STATUS.txt}, \texttt{K3\_SELECTION\_REPORT.md} --- checked
  2026-07-30, no hit beyond the control note itself and an identical
  regression-control comment in Stream 2's G0 checker). Any inbound brief
  presenting $s_{10}$'s $\U\oplus\lat{20}$ as a certified, stand-alone result
  should be treated as stale and returned for provenance, per the program's
  standing verification rule. \emph{[Update 2026-08-01: Stream~2 has since
  promoted the fiber-level Nikulin-complement computation for $s_{10}$ to a
  titled draft certificate (\texttt{G0\_NS\_genus\_cooper\_s10.json}, Tier~B,
  pending review), with its own negative-control suite. That certificate
  takes the $s_{10}$ transcendental-lattice identification as an \emph{input}
  and states explicitly that the reviewed monodromy-lattice derivation this
  item says is missing remains missing; the stance of this item --- no
  $s_{10}$ lattice claim is made in this paper --- is unchanged.]}
\end{itemize}

\subsection{Open questions}

\begin{enumerate}
\item Make Numerical Claim~\ref{nc:monodromy} exact. The natural route is
  exact hypergeometric connection formulae for $\Ltwo$, which would upgrade
  Propositions~\ref{prop:lattice} and~\ref{prop:usplit} to unconditional
  statements about the $s_7$ family and remove the only numerical input in
  this paper.
\item Close the rescaling ambiguity of Conjecture~\ref{conj:T} by a
  computation rather than by the framework's shape --- for instance by
  exhibiting a $\lambda>1$ obstruction intrinsic to the monodromy
  representation.
\item Prove Theorem~\ref{thm:obrien} in a proof assistant, discharging the one
  literature axiom of the formal development (\S\ref{sec:formalization}). This
  requires $q$-expansion machinery for weight-1 level-7 modular objects that
  the pinned library does not have.
\item Determine whether the $s_{18}$ partner exhibited in
  Corollary~\ref{cor:sporadic-partners} --- which appears not to have been
  written down before --- has any modular interpretation, or whether its
  non-integrality is an obstruction to one.
\item Establish the genus statement for the abstract lattice
  $\U\oplus\lat{14}$ independently. It is not needed for
  Proposition~\ref{prop:usplit}, which is constructive and therefore stronger
  than a genus argument for the specific lattice at hand, but it would be the
  natural framing for a classification statement.
\end{enumerate}

\subsection*{A further construction (placeholder, 2026-07-29)}

A further section, extending the lattice-theoretic results above toward an
explicit higher-dimensional construction, is planned but not yet drafted: an
in-progress feasibility check on the relevant base spaces has not concluded,
and drafting ahead of it risks stating a result that the check may
contradict. This placeholder will be replaced once that check resolves; no
content, method, or conclusion of the pending section is anticipated here.

% ----------------------------------------------------------------------------
% Generated-by: Opus 5 (Stream 1 paper agent), G0/G1 additions 2026-07-29 by
% Sonnet 5 (WP-P1) | Verified-by: framework statements vs paper/references.bib
% entries Dolgachev1996 / Doran2000 as recorded in PLAN.md row #19; rho=19 /
% T=3 conditionality and the (H1)/(H2) split vs Stream 2
% data/certificates/L3_IRREDUCIBLE.json step5_hodge_identification +
% conditional_on_step5; overlattice and rescaling ambiguities vs
% C2_cooper_s7_v4.json and briefs/STREAM2_U1_EXECUTION_2026_07_27.md §4.
% Claims #18-#20 of paper/PLAN.md §3; presented conditionally per PLAN.md §5.4.
% G0 (Prop.~g0complement, \S8.4): S2 briefs/G0_NS_GENUS_RESULT_2026_07_28.md
% (commits 9a386d9, 15f16c5) + decision log
% briefs/WP_S2G_X4_EXHIBITION_PLAN_2026_07_27.md §8 (commit 256017d, LIVE
% promotion + cross-model verification entry); independence caveats
% transcribed from the G0 brief's own "Two honest limits" paragraph, not
% invented here. G1 placeholder: PLAN.md scope ruling (pure mathematics, zero
% physics) kept in mind -- deliberately undescribed and unmotivated per WP-P1
% instructions; whether it belongs in this manuscript at all is flagged as a
% scope question for T0 in briefs/P1_DRAFT_STATUS_2026_07_29.md.
% U+<20> stub added 2026-07-30 by Sonnet 5 (T0 Directive 3): searched this
% repo (paper/PLAN.md row 17, 06-lattice.tex nc:s10-control) and Stream 2
% (~/SocrateAI-Scientific-Agora-K3-DarkMatter) PREDICTION.md, PROOF_STATUS.txt,
% K3_SELECTION_REPORT.md for a certified cooper_s10 monodromy-lattice
% derivation; none found -- only the discriminating-control number (already
% correctly tiered NC/(N)+(E) in this manuscript) and an identical
% regression-control comment in Stream 2's checkers/check_NS_genus_G0.py /
% data/certificates/G0_NS_genus_cooper_s7.json. Per Stream 2 Standing Rule 4,
% the directive's implied "certified cooper_s10 section" is not executed;
% this explicit stub is inserted instead. Deviation-with-cause recorded here
% and in the commit message.
% | Reviewed-by: pending T0 (Xavier)
% ----------------------------------------------------------------------------
